\RequirePackage{fix-cm}
\documentclass[smallextended,envcountsame]{svjour3}
\smartqed

\usepackage[T1]{fontenc}
\usepackage{amsmath,amssymb,mathtools}
\usepackage{booktabs}
\usepackage{microtype}
\usepackage[hidelinks]{hyperref}
\usepackage{url}

\newcommand{\R}{\mathbb R}
\newcommand{\ip}[2]{\langle #1,#2\rangle}
\newcommand{\norm}[1]{\lVert #1\rVert}

\title{A uniformly convex counterexample to global convergence of DFP
under weak Wolfe line search}
\titlerunning{A uniformly convex counterexample for DFP}

% SUBMISSION METADATA: replace the anonymous working-copy fields below before
% uploading the manuscript to Editorial Manager.  Do not infer these factual
% declarations from the local project owner name.
\author{Anonymous Author(s)}
\authorrunning{Anonymous Author(s)}
\institute{Author names, affiliations, corresponding-author email, and ORCID
will be inserted in the final submission file.}
\date{Received: date / Accepted: date}
\journalname{Mathematical Programming}

\begin{document}
\maketitle

\begin{abstract}
Does the classical Davidon--Fletcher--Powell (DFP) method converge globally
on every uniformly convex objective when its steps satisfy the standard weak
Wolfe conditions?  We give a negative answer.  We construct a twice
continuously differentiable function $f:\R^2\to\R$ with
$\frac12 I\preceq\nabla^2 f(x)\preceq\frac32 I$ for every $x$, a symmetric
positive definite initial inverse-Hessian approximation, and positive step
lengths satisfying weak Wolfe with $(c_1,c_2)=(1/4,3/4)$, but for which the
DFP gradient norms converge to a positive constant.  The construction begins
with an exact two-phase singular DFP orbit.  Along an invariant slow curve,
one metric eigenvalue tends to zero, the gradient amplitude remains positive,
and the iterates form an inward discrete spiral around a limiting circle.  We
prove a uniform all-pairs separation estimate for the spiral and use
pairwise-disjoint smooth corrections of a quadratic to interpolate every
endpoint gradient.  The corrections have Hessian norms tending to zero at
the circle, so their sum defines one global uniformly convex $C^2$ objective.
Every secant equation, DFP update, Armijo inequality, and weak-curvature
inequality then holds exactly.  Orthogonal direct sums extend the example to
every dimension $n\ge2$.
\keywords{Davidon--Fletcher--Powell method \and quasi-Newton method \and
weak Wolfe line search \and global convergence \and uniform convexity \and
counterexample}
\subclass{90C53 \and 65K05 \and 90C30}
\end{abstract}

\section{Introduction}
\label{sec:introduction}

The Davidon--Fletcher--Powell (DFP) update is one of the original variable
metric methods for unconstrained optimization
\cite{Davidon1991,FletcherPowell1963}.  Combined with a Wolfe line search
\cite{Wolfe1969}, it generates positive definite inverse-Hessian
approximations from positive-curvature secant pairs and is locally
superlinearly convergent under the familiar regularity assumptions.  Its
global behavior is subtler.  Powell proved convergence for uniformly convex
objectives under exact line search \cite{Powell1971}, but exact search supplies
the orthogonality $g_{k+1}^{T}s_k=0$, an identity that an ordinary weak Wolfe
line search does not retain.

The endpoint nature of the difficulty is already visible in the convex
Broyden class.  Byrd, Nocedal, and Yuan established a global self-correction
mechanism for a restricted class but explicitly excluded the DFP endpoint
\cite{ByrdNocedalYuan1987}.  Yuan later proved convergence under additional
trajectory hypotheses, including eventual monotonicity of the gradient norm
or summability of the step lengths \cite{Yuan1995}.  Xu obtained another
conditional theorem involving an eventual monotonicity condition and the
parameter restriction $c_2<(m/M)^3$ when
$mI\preceq\nabla^2f\preceq MI$ \cite{Xu1997}.  These results left open whether
uniform convexity and the standard weak Wolfe conditions alone suffice.  The
question was highlighted in early overviews by Nocedal and Fletcher
\cite{Nocedal1992,Fletcher1994}.  A paper published in 2026 still states this
precise question as open and changes the algorithm by adding a
projection/correction mechanism
\cite{YuanZhouPham2026}.

We show that the missing hypothesis cannot be removed: classical,
unmodified DFP can fail globally even in dimension two.  The obstruction is
not a finite ill-conditioning transient.  We construct an infinite orbit on
which a small metric eigenvalue decays through a parabolic two-phase drift,
while the gradient amplitude converges to a nonzero limit.  The associated
physical frame rotates by a nonsummable amount, and the iterates wind inward
toward a circle.  The curvature perturbations used to discover the orbit tend
to the identity, which is essential for realizing the entire sequence by the
gradient of one $C^2$ function near its accumulation set.

The main contributions are as follows.
\begin{enumerate}
\item We derive an exact two-leg DFP map, remove its apparent singularity, and
      identify a one-dimensional $C^7$ invariant slow curve.  On that curve,
      the singular scale obeys
      $\epsilon_{j+1}=\epsilon_j-\frac32\epsilon_j^4+O(\epsilon_j^5)$.
\item We compute the physical drift of this orbit.  The gradient amplitude
      loses only a summable amount, whereas the frame rotation is
      nonsummable.  This creates an inward discrete spiral with a full circle
      as its accumulation set.
\item We prove a uniform separation estimate for every pair of endpoints,
      including near returns after one or many turns.  Pairwise-disjoint
      cutoff corrections then realize all endpoint jets and extend with a
      zero second-order jet on the limiting circle.
\item We obtain one global $C^2$ function satisfying
      $\frac12I\preceq\nabla^2f\preceq\frac32I$, and verify the original DFP
      recurrence and both weak Wolfe inequalities exactly.  The gradient
      norms converge to a positive constant, giving a counterexample in every
      dimension $n\ge2$.
\end{enumerate}

The example is consistent with all of the conditional results just
mentioned.  Its gradient norm decreases on the first leg of each cycle and
increases on the second; its step lengths in state space are asymptotic to
$2\epsilon_j^2$ and $\epsilon_j^2$, whose sum diverges; and the selected value
$c_2=3/4$ is larger than $(m/M)^3=1/27$ for
$(m,M)=(1/2,3/2)$.  Section~\ref{sec:verification} gives the precise
comparison.

The paper is organized as follows.  Section~\ref{sec:statement} states the
counterexample and records the exact one-step algebra.  The singular
two-phase orbit is constructed in Section~\ref{sec:two-phase}.
Section~\ref{sec:spiral} proves its asymptotics and uniform separation.
Section~\ref{sec:realization} constructs the global objective, and
Section~\ref{sec:verification} completes the Wolfe, convergence-failure, and
known-results checks.  Appendix~\ref{app:algebra} exposes the algebra behind
the decisive two-leg coefficients.

\section{The DFP iteration and the counterexample}
\label{sec:statement}

For $g_k=\nabla f(x_k)$, classical inverse-form DFP is
\begin{align*}
d_k&=-H_kg_k,&s_k&=\alpha_kd_k,&x_{k+1}&=x_k+s_k,\\
y_k&=g_{k+1}-g_k,&
H_{k+1}&=H_k-
\frac{H_ky_ky_k^TH_k}{y_k^TH_ky_k}
+\frac{s_ks_k^T}{s_k^Ty_k}.
\end{align*}
The standard weak Wolfe conditions are
\begin{align}
f(x_{k+1})&\le f(x_k)+c_1g_k^Ts_k,\label{eq:armijo}\\
g_{k+1}^Ts_k&\ge c_2g_k^Ts_k,
\label{eq:weak-curvature}
\end{align}
where $0<c_1<c_2<1$.

\begin{theorem}[Main theorem]\label{thm:main}
There exist $f\in C^2(\R^2)$, $x_0\in\R^2$, $H_0\succ0$, and positive
step lengths $\alpha_k$ such that
\[
\frac12I\preceq\nabla^2f(x)\preceq\frac32I
\qquad(x\in\R^2),
\]
all classical DFP iterates are well defined, and every step satisfies
\eqref{eq:armijo}--\eqref{eq:weak-curvature} with
\[
c_1=\frac14,\qquad c_2=\frac34,
\]
but
\[
\|\nabla f(x_k)\|\longrightarrow G_\infty>0.
\]
The same conclusion holds in every dimension $n\ge2$ by an orthogonal direct
sum.
\end{theorem}

The auxiliary matrices used below generate an abstract sequence of secants.
They are not asserted to be the Hessians of the final objective.  The final
objective interpolates the endpoint gradients exactly, so its actual averaged
Hessians produce the same secants automatically.

\section{The exact two-phase DFP map}
\label{sec:two-phase}

At a cycle boundary, use an oriented orthogonal frame $R=(e,e^\perp)$ and
write
\begin{equation}\label{eq:canonical-state}
H=R\begin{pmatrix}hpr^2&0\\0&h\end{pmatrix}R^T,
\qquad
g=GR\binom{1}{pr},
\end{equation}
where $r,p,h,G>0$.  For a symmetric positive definite matrix $A$ and a
number $\tau>0$, define
\begin{equation}\label{eq:abstract-step}
\alpha=\tau\frac{g^THg}{(Hg)^TA(Hg)},\qquad
s=-\alpha Hg,\qquad y=As,\qquad g_+=g+y,
\end{equation}
and update $H$ by DFP.  Then, with
\[
q=-g^Ts>0,\qquad t=s^Ty,
\]
one has the exact identity
\begin{equation}\label{eq:line-ratio}
\frac tq=\tau.
\end{equation}
Indeed, $q=\alpha g^THg$ and
$t=\alpha^2(Hg)^TA(Hg)=\tau\alpha g^THg=\tau q$.

Throughout the proof, if $e=(e_1,e_2)^T$, then
$e^\perp=(-e_2,e_1)^T$, so $R$ is positively oriented.  Vector norms are
Euclidean and matrix norms are spectral norms.

The following formula is the algebraic starting point for every expansion in
the paper.  It also separates the exact DFP calculation from the later
asymptotic choices.

\begin{proposition}[Exact one-step map]\label{prop:one-step}
Let $v=Hg$, $\delta=g^THg$, $w=Av$, $\beta=v^TAv$, and
$\gamma=w^THw$.  The step in \eqref{eq:abstract-step} satisfies
\begin{align}
g_+&=g-\tau\frac{\delta}{\beta}w,
\label{eq:one-step-g}\\
H_+&=H-\frac{Hww^TH}{\gamma}+\frac{vv^T}{\beta}.
\label{eq:one-step-H}
\end{align}
In particular, the matrix update is independent of the scale parameter
$\tau$, and all denominators in \eqref{eq:one-step-g}--\eqref{eq:one-step-H}
are positive.
\end{proposition}

\begin{proof}
Equation \eqref{eq:abstract-step} gives
$s=-\alpha v$ and $y=-\alpha w$.  Therefore
$s^Ty=\alpha^2\beta$, $y^THy=\alpha^2\gamma$, and the two factors
$\alpha^2$ cancel in the rank-one DFP terms.  Substitution gives
\eqref{eq:one-step-H}, while $g_+=g+y$ gives
\eqref{eq:one-step-g}.  Positivity follows from $H\succ0$, $A\succ0$, and
$v\ne0$.
\end{proof}

Set $r=\epsilon^2$.  During a complete cycle use, in the current oriented
eigenframe, the controls
\begin{equation}\label{eq:two-controls}
A_1(\epsilon)=
\begin{pmatrix}1&\epsilon\\ \epsilon&1\end{pmatrix},
\quad \tau_1=\frac23;
\qquad
A_2(\epsilon)=
\begin{pmatrix}1&-2\epsilon\\ -2\epsilon&1\end{pmatrix},
\quad \tau_2=\frac13.
\end{equation}
We restrict the initial scale so that
$0<\epsilon\le\epsilon_0<1/4$; hence both controls are positive definite
(indeed, their spectra lie in $[1/2,3/2]$).
The same $\epsilon$ is used on both legs.  Between the legs and after the
second leg, diagonalize the new $H$, orient its low eigenvector to have
positive inner product with the new gradient, and use the corresponding
coordinates in \eqref{eq:canonical-state}.

We use the following finite-smooth invariant-graph statement.  It is included
to make explicit both the regularity needed below and the fact that the full
map need not be invertible.

\begin{lemma}[Local invariant graph for a map]\label{lem:graph-transform}
Let $\nu\ge2$ and let $F=(F_c,F_s)$ be a $C^\nu$ map from a neighborhood of
$(0,0)\in\R\times\R^m$ into $\R\times\R^m$.  Suppose
\[
F(0,0)=(0,0),\qquad
DF(0,0)=\begin{pmatrix}1&0\\0&L\end{pmatrix},
\qquad \rho(L)<1.
\]
Then, after restricting the neighborhood, there is a $C^\nu$ function
$z=\zeta(u)$ with $\zeta(0)=D\zeta(0)=0$ whose graph is locally forward
invariant.  The conclusion does not require $L$, or the full derivative
$DF(0,0)$, to be invertible.
\end{lemma}

\begin{proof}
Choose an equivalent norm in the stable space for which
$\norm{L}<\lambda<1$.  For every graph with sufficiently small Lipschitz
constant, the center projection
$u\mapsto F_c(u,\zeta(u))$ has derivative uniformly close to one and is
therefore a local $C^\nu$ diffeomorphism.  Its inverse defines the graph
transform
\[
(\mathcal T\zeta)(\bar u)
=F_s(u,\zeta(u)),
\qquad \bar u=F_c(u,\zeta(u)).
\]
After the usual local cutoff, the stable contraction and the small nonlinear
derivatives make $\mathcal T$ a contraction on a closed small Lipschitz cone.
Its fixed graph is forward invariant.  For the prescribed finite $\nu$, the
neighborhood can also be reduced so that the stable contraction multiplied
by the $\nu$th power of the inverse center-derivative bound is still less than
one; this is possible because these two factors tend to $\norm L<1$ and $1$.
The differentiated graph transforms are then contractions through order
$\nu$, giving $C^\nu$ regularity.  This is the finite-smooth graph-transform form of
the discrete center-manifold theorem; see, for example,
Karydas and Schinas \cite{KarydasSchinas1992}.  The construction only inverts
the one-dimensional center projection, so a zero eigenvalue of $L$ causes no
obstruction.
\end{proof}

\begin{lemma}[Blown-up map and invariant slow curve]
\label{lem:slow-curve}
The exact two-leg map in the variables $(\epsilon,p,h)$ has a removable
real-analytic extension near $(0,2,1)$.  Its derivative there has one center
eigenvalue $1$ and transverse eigenvalues $-1/9$ and $0$.  It therefore has a
locally invariant $C^7$ center curve
\[
p=p(\epsilon),\qquad h=h(\epsilon),
\]
on which
\begin{align}
p(\epsilon)&=2+\frac{198}{5}\epsilon^3
-\frac95\epsilon^4+O(\epsilon^5),
&h(\epsilon)&=1+8\epsilon^3+O(\epsilon^5),
\label{eq:center-curve}\\
\epsilon_+&=\epsilon-\frac32\epsilon^4
+\frac54\epsilon^5+O(\epsilon^6).
\label{eq:epsilon-map}
\end{align}
\end{lemma}

\begin{proof}
All operations in one DFP step are rational in the matrix entries and the
secant pair.  At the blown-up base point, every denominator left after
cancelling the explicit powers of $\epsilon$ is positive.  The two eigenvalues
of the limiting matrix are $0$ and $1$, hence its oriented spectral frame is
analytic.  Direct cancellation therefore gives the claimed analytic
extension.  If $\widehat r_+$ denotes the canonical radius after two legs,
then $\widehat r_+/\epsilon^2$ has a positive analytic extension with value
one.  We use the signed branch
\[
\epsilon_+=\epsilon
\sqrt{\widehat r_+/\epsilon^2}.
\]
It agrees with the positive square root when $\epsilon>0$ and makes the
three-variable map analytic across $\epsilon=0$.

Here is the explicit removable-factor check.  After leg $i\in\{1,2\}$,
write the state in its new oriented frame as
\[
\lambda_-^{(i)}=\epsilon^4L_i,\qquad
\lambda_+^{(i)}=\mathcal H_i,\qquad
R_i^Tg_i
=G\binom{\mathcal G_i}{\epsilon^2U_i},\qquad
r_i=\epsilon^2\mathcal R_i.
\]
Direct substitution in \eqref{eq:abstract-step} and the DFP determinant
identity
\[
\det H_+=\det H\,\frac{s^Ty}{y^THy}
\]
show that all displayed factors are analytic.  Their base-point
values are
\[
\begin{array}{c|cccccc}
i&L_i&\mathcal H_i&\mathcal G_i&U_i&\mathcal R_i&p_i\\ \hline
1&2&1&1&1&2&1/2\\
2&2&1&1&2&1&2.
\end{array}
\]
Indeed,
\[
r_i=\frac{\lambda_-^{(i)}}
{\lambda_+^{(i)}
 ((R_i^Tg_i)_2/(R_i^Tg_i)_1)}
=\epsilon^2\frac{L_i\mathcal G_i}{\mathcal H_iU_i}.
\]
On each leg the four quantities
$g^THg$, $(Hg)^TA(Hg)$, $s^TAs$, and $(As)^TH(As)$ equal
$\epsilon^4$ times analytic factors with positive base-point values.
For the first leg these four limiting factors are $(6,4,4,4)$; for the
second they are $(3,1,1,1)$.  For example, before the first update,
\[
g^THg=hp(p+1)\epsilon^4,
\qquad
(Hg)^TA_1(Hg)=h^2p^2\epsilon^4
 (1+2\epsilon^3+\epsilon^4).
\]
Consequently every step length, DFP denominator, spectral coordinate, and
canonical-state quotient has a removable analytic extension.  In particular
$\widehat r_+/\epsilon^2=\mathcal R_2\to1$, as used in the signed branch.

The cancellation also gives an explicit certificate for the limiting map.
If $p$ and $h$ are left independent while $\epsilon\to0$, then
\begin{align}
\left.\frac{r_+}{\epsilon^2}\right|_{\epsilon=0}
&=\frac{9hp(p+1)}{2\{9hp+(p+1)^2\}},
\label{eq:limiting-radius}\\
\left.p_+\right|_{\epsilon=0}
&=\frac{4\{9hp+(p+1)^2\}^2}
{81hp(p+1)^2},
&\left.h_+\right|_{\epsilon=0}&=1.
\label{eq:limiting-ph}
\end{align}
At $(p,h)=(2,1)$, the radius factor equals one and the limiting map fixes
$(2,1)$.  Differentiating \eqref{eq:limiting-radius}--\eqref{eq:limiting-ph}
yields
\begin{equation}\label{eq:limiting-jacobian}
D(\epsilon_+,p_+,h_+)(0,2,1)
=\begin{pmatrix}
1&0&0\\
0&-1/9&2/3\\
0&0&0
\end{pmatrix}.
\end{equation}
Thus the center and transverse spectra claimed in the lemma are visible
directly from the exact limiting formulas.

For completeness, introduce independent small variables $r$ and $b$, and put
\[
p=2+Pb r,\qquad h=1+Jb r.
\]
Substitution in the two exact DFP updates, followed by the quadratic formula
for the two simple eigenvalues, gives the following formulas.  Here
$C=x-g$, $e$ is the physical low eigenvector, and $\phi$ is its unwrapped
angle:
\begin{align*}
r_+-r
 &=\frac{b(6J+5P-300)}{18}r^2+O(r^3),\\
p_+-2
 &=\frac{b(6J-P+348)}9r+O(r^2),\\
h_+-1&=8br+O(r^2),\\
\frac{G_+}{G}-1
 &=\frac{b^2(24J-4P+384)-117}{18}r^2+O(r^3),\\
\phi_+-\phi&=-3r+O(r^2),\\
e^T(C_+-C)
 &=-\frac{2b^2(6J-P+96)}9Gr^2+O(G|b|r^3).
\end{align*}
The remainders are uniform for bounded $(b,P,J)$ in a fixed neighborhood.
In particular, the leading transverse map is
\begin{equation}\label{eq:transverse-map}
P_+=\frac{6J-P+348}{9}+o(1),
\qquad J_+=8+o(1).
\end{equation}
Thus its linear part is
\[
L=\begin{pmatrix}-1/9&2/3\\0&0\end{pmatrix},
\]
and the fixed coefficient is
\[
(P,J)=\left(\frac{198}{5},8\right).
\]
Lemma~\ref{lem:graph-transform} applies with $\nu=7$: the center eigenvalue is
$1$, while the two transverse eigenvalues $-1/9$ and $0$ lie strictly inside
the unit disk.  The analyticity proved above supplies all required
derivatives, and the zero eigenvalue belongs to the stable block.  We obtain a
$C^7$ invariant graph.  Coefficient matching can be displayed explicitly.
Put
\[
p=2+P_3\epsilon^3+P_4\epsilon^4+O(\epsilon^5),
\qquad
h=1+H_3\epsilon^3+H_4\epsilon^4+O(\epsilon^5).
\]
The $p$- and $h$-components of the graph-invariance equation give
\begin{equation}\label{eq:graph-coefficient-system}
3H_3-5P_3+174=0,\qquad 8-H_3=0,\qquad
3H_4-5P_4-9=0,\qquad H_4=0.
\end{equation}
Hence
\[
(P_3,H_3,P_4,H_4)=\left(\frac{198}{5},8,-\frac95,0\right).
\]
For the signed center coordinate, the next two coefficients of
$\epsilon_+-\epsilon$ are
\[
\frac{6H_3+5P_3-300}{36}=-\frac32,
\qquad
\frac{6H_4+5P_4+54}{36}=\frac54.
\]
This proves \eqref{eq:center-curve}--\eqref{eq:epsilon-map}.  In particular,
the leading relation is also recovered from the displayed bivariate formula:
$r_+=r-3\epsilon r^2+O(\epsilon^6)$ when
$b=\epsilon$ and $r=\epsilon^2$.
Finally, \eqref{eq:epsilon-map} gives
$0<\epsilon_+<\epsilon$ for every sufficiently small $\epsilon>0$.
Since $p(0)=2$, $h(0)=1$, and the center curve is continuous, its entire
forward orbit then remains in the positive region $p,h,\epsilon>0$.
\end{proof}

Fix first a uniform Taylor neighborhood
$0<\epsilon\le\bar\epsilon<1/4$ on which all analytic denominators and the
invariant graph above are defined.  All constants implicit below are chosen
on this fixed neighborhood and are independent of the final initial scale
$\epsilon_0\le\bar\epsilon$.  Choose such an $\epsilon_0>0$, take
$(p_0,h_0)=(p(\epsilon_0),h(\epsilon_0))$, $G_0=1$, and $R_0=I$ in
\eqref{eq:canonical-state}, and set $x_0=g_0$, so that $C_0=x_0-g_0=0$.
The preceding construction now defines an exact infinite abstract DFP orbit.

\begin{lemma}[Physical drift on the slow curve]\label{lem:physical-drift}
Let $\phi_j$ be the unwrapped angle of the low eigenvector at cycle boundary
$j$, let $G_j$ be the amplitude in \eqref{eq:canonical-state}, and define
\[
C_k=x_k-g_k.
\]
Along the invariant curve,
\begin{align}
\frac{G_{j+1}}{G_j}
&=1-\frac{13}{2}\epsilon_j^4+O(\epsilon_j^6),
\label{eq:G-map}\\
\phi_{j+1}-\phi_j
&=-3\epsilon_j^2+O(\epsilon_j^4),
\label{eq:phi-map}\\
C_{2j+2}-C_{2j}
&=-\frac{116}{5}G_j\epsilon_j^6e_j
+O(G_j\epsilon_j^7),
\label{eq:C-map}\\
\|C_{2j+1}-C_{2j}\|&=O(G_j\epsilon_j^3).
\label{eq:C-half}
\end{align}
The two endpoint polar-angle increments within a cycle are respectively
\begin{equation}\label{eq:half-angles}
-2\epsilon_j^2+o(\epsilon_j^2),
\qquad -\epsilon_j^2+o(\epsilon_j^2).
\end{equation}
Here $e_j$ is the physical low unit eigenvector.
\end{lemma}

\begin{proof}
After inserting the leading slow-graph coefficients
$(P,J)=(198/5,8)$, the same exact bivariate expansion used in
\eqref{eq:transverse-map} gives, before putting $b=\epsilon$,
\begin{align*}
\frac{G_+}{G}
&=1+\frac{232b^2-65}{10}r^2+O(r^3),\\
\phi_+-\phi&=-3r+O(r^2),\\
e^T(C_+-C)&=-\frac{116}{5}Gb^2r^2+O(G|b|r^3).
\end{align*}
These remainders are uniform for bounded $b$ in the fixed Taylor
neighborhood.  The last remainder is deliberately not written as
$o(Gb^2r^2)$ for two independent variables.  Along the path used here,
$b=\epsilon$ and $r=\epsilon^2$, it is
$O(G\epsilon^7)=o(G\epsilon^6)$.
The four center-curve coefficients in \eqref{eq:center-curve}, substituted
into the exact analytic map, sharpen these relations to
\begin{align*}
\frac{G_+}{G}
&=1-\frac{13}{2}\epsilon^4
+\frac{116}{5}\epsilon^6+O(\epsilon^7),\\
\phi_+-\phi&=-3\epsilon^2+O(\epsilon^4),\\
R_j^T(C_{2j+2}-C_{2j})/G_j
&=\binom{-\frac{116}{5}\epsilon^6+O(\epsilon^7)}
{O(\epsilon^8)}.
\end{align*}
Along the single-parameter path $b=\epsilon$, $r=\epsilon^2$, the first-leg
formula is
\[
C_{2j+1}-C_{2j}
=G_jR_j\binom{2\epsilon^3+o(\epsilon^3)}
{2\epsilon^5+o(\epsilon^5)}.
\]
These relations prove
\eqref{eq:G-map}--\eqref{eq:C-half}.  Direct expansion of the two oriented
eigenframes and gradients gives \eqref{eq:half-angles}.
\end{proof}

\section{Asymptotics and the discrete spiral}
\label{sec:spiral}

\begin{lemma}[Scalar asymptotics]\label{lem:scalar-asymptotics}
There are limits $G_\infty>0$ and $C_\infty\in\R^2$ such that
\begin{align}
\epsilon_j&\sim\left(\frac92j\right)^{-1/3},
\label{eq:epsilon-asymptotic}\\
G_j-G_\infty&\sim\frac{13}{3}G_\infty\epsilon_j,
\label{eq:G-tail}\\
\|C_{2j}-C_\infty\|&=O(\epsilon_j^3),
\qquad
\|C_{2j+1}-C_\infty\|=O(\epsilon_j^3).
\label{eq:C-tail}
\end{align}
Moreover $\phi_j\to-\infty$, and the accumulation set of the endpoints is
the full circle
\begin{equation}\label{eq:limit-circle}
\Gamma=\{C_\infty+G_\infty v:\|v\|=1\}.
\end{equation}
\end{lemma}

\begin{proof}
The standard reciprocal-power argument applied to
\eqref{eq:epsilon-map} gives \eqref{eq:epsilon-asymptotic}.  In particular,
\[
\sum_j\epsilon_j^4<\infty,
\qquad
\sum_j\epsilon_j^2=\infty.
\]
The positive product in \eqref{eq:G-map} therefore converges to
$G_\infty>0$.  Comparing the tails of \eqref{eq:G-map} with
\eqref{eq:epsilon-map} gives
\[
\sum_{\ell\ge j}\epsilon_\ell^4\sim\frac23\epsilon_j,
\]
and hence \eqref{eq:G-tail}.  The series in \eqref{eq:C-map} is absolutely
convergent, with
\[
\sum_{\ell\ge j}\epsilon_\ell^6=O(\epsilon_j^3).
\]
Combining this with \eqref{eq:C-half} proves \eqref{eq:C-tail}.

Equation \eqref{eq:phi-map} and the divergent sum of $\epsilon_j^2$ imply
$\phi_j\to-\infty$.  The angular increments tend to zero, so every angle
modulo $2\pi$ is approached.  Equations \eqref{eq:G-tail} and
\eqref{eq:C-tail} then give \eqref{eq:limit-circle}.
\end{proof}

For an endpoint $x_k$ in cycle $j$, write $r(k)=\epsilon_j^2$.

\begin{lemma}[Uniform separation]\label{lem:separation}
After decreasing $\epsilon_0$ if necessary, there is a constant $c_*>0$ such
that, for every $k$,
\begin{equation}\label{eq:separation}
\operatorname{dist}\!\left(
x_k,(\{x_\ell:\ell\ge0\}\cup\Gamma)\setminus\{x_k\}
\right)
\ge c_*r(k).
\end{equation}
\end{lemma}

\begin{proof}
At a cycle boundary, complex notation in the oriented physical frame gives
\begin{equation}\label{eq:polar-boundary}
x_{2j}-C_\infty
=G_je^{i\phi_j}\bigl(1+2i\epsilon_j^2+o(\epsilon_j^2)\bigr).
\end{equation}
The endpoint angles are strictly decreasing for small $\epsilon_0$.  By
\eqref{eq:half-angles}, there are constants
$0<c_\theta<C_\theta$, fixed on the preliminary neighborhood, such that each
consecutive endpoint gap lies between $c_\theta r_j$ and
$C_\theta r_j$.  Enlarging $C_\theta$ if necessary, any block of $N$ cycles
whose scales are comparable rotates by at most $NC_\theta r_j$.  During one
complete turn, $\epsilon$ changes by only a relative $O(\epsilon_j)$.

For both phases $\sigma\in\{0,1\}$, the polar radius satisfies
\begin{equation}\label{eq:phase-radius}
\|x_{2j+\sigma}-C_\infty\|
=G_j+O(\epsilon_j^3)=G_j+o(r_j),
\qquad r_j=\epsilon_j^2.
\end{equation}
More quantitatively, there is a modulus $\omega(\eta)\downarrow0$ such that,
whenever $\epsilon_0\le\eta$,
\begin{equation}\label{eq:phase-radius-uniform}
\left|\norm{x_{2j+\sigma}-C_\infty}-G_j\right|
\le \omega(\eta)r_j
\qquad(\sigma=0,1;\ j\ge0).
\end{equation}
Fix two distinct endpoints, with the first in cycle $j$ and the second in
cycle $\ell\ge j$, and use the earlier scale $r_j$.  The following three
cases are exhaustive.

Choose once and for all
\begin{equation}\label{eq:kappa-choice}
\frac1{\sqrt2}<\kappa<1.
\end{equation}
First, suppose $\epsilon_\ell\le\kappa\epsilon_j$.
Equation \eqref{eq:G-tail} and
\eqref{eq:phase-radius} give
\[
G_j-G_\ell
=\frac{13}{3}G_\infty(\epsilon_j-\epsilon_\ell)(1+o(1))
\ge c\epsilon_j\gg r_j,
\]
so the radial component alone proves the required separation.

Next suppose the two scales are comparable and their circular angular
separation is at least $r_j/4$.  Both polar radii have a uniform positive
lower bound, and the chord formula gives distance at least $c r_j$.
This case includes adjacent endpoints because their angular gaps are bounded
below by a fixed multiple of $r_j$.

It remains to consider comparable scales whose unwrapped angular difference
is $2\pi m+\zeta$, where $|\zeta|<r_j/4$.  The lower bound on consecutive
angular gaps excludes $m=0$: after decreasing $\epsilon_0$, each consecutive
gap is at least $r_t/2$, while comparability and
\eqref{eq:kappa-choice} give
$r_t/2\ge\kappa^2r_j/2>r_j/4$.  Thus $m\ge1$.

Let $N=\ell-j$.  Here $N\ge1$; by enlarging $C_\theta$ once, the two endpoint
phase offsets are absorbed into the same $NC_\theta r_j$ bound.  While the
scales are comparable, the amplitude expansion
\eqref{eq:G-map} provides a fixed $c_G>0$ such that
$G_t-G_{t+1}\ge c_G r_j^2$.  The angular displacement gives the explicit
cycle-count inequality
\begin{equation}\label{eq:cycle-count}
NC_\theta r_j\ge 2\pi m-r_j/4,
\end{equation}
and consequently
\begin{equation}\label{eq:radial-return-gap}
G_j-G_\ell
\ge Nc_G r_j^2
\ge \frac{c_G}{C_\theta}(2\pi m-r_j/4)r_j
\ge c_0m r_j
\end{equation}
for a uniform $c_0>0$.  We now choose the final upper bound on
$\epsilon_0$ so that $2\omega(\epsilon_0)\le c_0/2$ in
\eqref{eq:phase-radius-uniform}.  Thus the two phase-radius errors are less
than half of the main radial gap in \eqref{eq:radial-return-gap}, uniformly in
the pair and in $m$.  For $m=1$, summing the leading terms recovers the
sharper one-turn
formula
\begin{equation}\label{eq:one-turn-drop}
G_j-G_\ell=\frac{13\pi}{3}G_jr_j(1+o(1)).
\end{equation}
This proves the estimate for endpoints in different phases, on neighboring
turns, and on multiple turns.  It does not assume that the physical endpoint
radius is monotone at every individual cycle.

Finally, \eqref{eq:G-tail}, \eqref{eq:C-tail}, and
\eqref{eq:polar-boundary} imply
\[
\operatorname{dist}(x_k,\Gamma)
=\frac{13}{3}G_\infty\epsilon_j(1+o(1))\gg r(k).
\]
All estimates are uniform on the invariant curve after taking one sufficiently
small $\epsilon_0$.  This proves \eqref{eq:separation}.
\end{proof}

\section{A global \texorpdfstring{$C^2$}{C2} uniformly convex realization}
\label{sec:realization}

Define the closed set
\[
\mathcal E=\Gamma\cup\{x_k:k\ge0\}
\]
and put
\begin{equation}\label{eq:dk-rhok}
d_k=\operatorname{dist}(x_k,\mathcal E\setminus\{x_k\}),
\qquad \rho_k=\frac14d_k.
\end{equation}
By Lemma~\ref{lem:separation}, $\rho_k\ge c r(k)$; the adjacent endpoint
also gives $\rho_k\le Cr(k)$.  The balls $B(x_k,\rho_k)$ are pairwise
disjoint and avoid $\Gamma$.

Let
\begin{equation}\label{eq:correction-vector}
a_k=g_k-(x_k-C_\infty)=C_\infty-C_k.
\end{equation}
Lemma~\ref{lem:scalar-asymptotics} gives
\begin{equation}\label{eq:ak-bound}
\|a_k\|=O(\epsilon_j^3)=O(\epsilon_jr(k)).
\end{equation}
Choose one $\chi\in C_c^\infty(B(0,1))$ with $\chi\equiv1$ on
$B(0,1/3)$, and define
\begin{equation}\label{eq:bumps}
\psi_k(z)=
\chi\!\left(\frac{z-x_k}{\rho_k}\right)a_k^T(z-x_k),
\qquad
\Psi(z)=\sum_{k=0}^\infty\psi_k(z).
\end{equation}

\begin{lemma}[Disjoint-bump extension]\label{lem:bump-extension}
The function $\Psi$, extended by zero on $\Gamma$, belongs to
$C^2(\R^2)$.  It satisfies
\begin{equation}\label{eq:bump-jets}
\Psi(x_k)=0,\qquad \nabla\Psi(x_k)=a_k,
\end{equation}
and
\begin{equation}\label{eq:bump-hessian}
\sup_{z\in\R^2}\|\nabla^2\Psi(z)\|
\le K\epsilon_0.
\end{equation}
\end{lemma}

\begin{proof}
Write $u=z-x_k$ and $\xi=u/\rho_k$.  Direct differentiation gives
\begin{align*}
\nabla\psi_k(z)
&=\chi(\xi)a_k
+\rho_k^{-1}(a_k^Tu)\nabla\chi(\xi),\\
\nabla^2\psi_k(z)
&=\rho_k^{-1}\{a_k\nabla\chi(\xi)^T
+\nabla\chi(\xi)a_k^T\}
+\rho_k^{-2}(a_k^Tu)\nabla^2\chi(\xi).
\end{align*}
On the support, $\|u\|\le\rho_k$.  Hence the fixed cutoff together with
\eqref{eq:dk-rhok} and \eqref{eq:ak-bound} implies
\begin{align*}
\|\psi_k\|_\infty&=O(\epsilon_jr(k)^2),\\
\|\nabla\psi_k\|_\infty&=O(\epsilon_jr(k)),\\
\|\nabla^2\psi_k\|_\infty&=O(\epsilon_j).
\end{align*}
At most one summand is nonzero at any point.  Away from $\Gamma$, the supports
are locally finite.  Along any sequence approaching $\Gamma$ through the
supports, the displayed function, gradient, and Hessian bounds all tend to
zero.  More explicitly, \eqref{eq:G-tail} and \eqref{eq:C-tail} give
\[
\operatorname{dist}(x_k,\Gamma)=\Theta(\epsilon_j),
\qquad \rho_k=O(\epsilon_j^2).
\]
Hence every point $z$ in the $k$th support has
$\operatorname{dist}(z,\Gamma)=\Theta(\epsilon_j)$, and
\[
\frac{|\psi_k(z)|}{\operatorname{dist}(z,\Gamma)^2}=O(\epsilon_j^3),
\qquad
\frac{\|\nabla\psi_k(z)\|}{\operatorname{dist}(z,\Gamma)}
=O(\epsilon_j^2),
\qquad
\|\nabla^2\psi_k(z)\|=O(\epsilon_j).
\]
These second-order difference quotients show that the zero value, gradient,
and Hessian on $\Gamma$ are the genuine limiting jets.  Thus $\Psi\in C^2$.
The cutoff is constant near each center, which gives
\eqref{eq:bump-jets}.  The constants in the separation estimate, the upper
bound $\rho_k\le Cr(k)$, and the derivatives of the fixed cutoff were chosen
on the preliminary neighborhood $\epsilon\le\bar\epsilon$.  Hence they
produce one constant $K$, independent of the later choice of
$\epsilon_0\le\bar\epsilon$.  Disjointness and
$\epsilon_j\le\epsilon_0$ then give \eqref{eq:bump-hessian}.
\end{proof}

Now set
\begin{equation}\label{eq:objective}
f(z)=\frac12\|z-C_\infty\|^2+\Psi(z).
\end{equation}
After $K$ has been fixed, take
$\epsilon_0\le\min\{\bar\epsilon,(2K)^{-1}\}$, together with the earlier
smallness requirements.  Then the right side of
\eqref{eq:bump-hessian} is at most $1/2$, which gives the global Hessian
bounds in Theorem~\ref{thm:main}.  At every endpoint,
\begin{equation}\label{eq:endpoint-data}
f(x_k)=\frac12\|x_k-C_\infty\|^2,
\qquad
\nabla f(x_k)=x_k-C_\infty+a_k=g_k.
\end{equation}
Thus the abstract sequence is an exact gradient sequence of one $C^2$
uniformly convex objective.  In particular, its $y_k$ are the actual gradient
secants, and every previously defined DFP update is unchanged.

\section{Completion and verification of the counterexample}
\label{sec:verification}

\subsection{Armijo decrease and weak curvature}

Let $q_k=-g_k^Ts_k>0$.  Because every bump vanishes at every endpoint,
\eqref{eq:endpoint-data} gives
\begin{align}
\frac{f(x_k)-f(x_{k+1})}{q_k}
&=1-\frac{\|s_k\|^2}{2q_k}
+\frac{a_k^Ts_k}{q_k}.
\label{eq:armijo-ratio}
\end{align}
The local state gives $\|s_k\|=\Theta(r(k))$ and
$q_k=\Theta(r(k)^2)$.  Equations \eqref{eq:ak-bound} and
\eqref{eq:line-ratio}, together with $A_i(\epsilon)\to I$, imply
\[
\frac{a_k^Ts_k}{q_k}=O(\epsilon_j),
\qquad
\frac{\|s_k\|^2}{q_k}=\tau_k+O(\epsilon_j).
\]
Therefore the two limits of \eqref{eq:armijo-ratio} are
\[
1-\frac12\frac23=\frac23,
\qquad
1-\frac12\frac13=\frac56.
\]
After one final reduction of $\epsilon_0$, every ratio is at least $1/2$;
Armijo holds with $c_1=1/4$.

Weak curvature is exact and value-independent:
\[
g_{k+1}^Ts_k=g_k^Ts_k+s_k^Ty_k
=-(1-\tau_k)q_k.
\]
Both $\tau_k\in\{2/3,1/3\}$ satisfy
$\tau_k\ge1-c_2=1/4$, proving \eqref{eq:weak-curvature} for $c_2=3/4$.

\subsection{Failure of convergence and extension to higher dimensions}

Finally, Lemma~\ref{lem:scalar-asymptotics} and
\eqref{eq:canonical-state} give the asserted limit at cycle boundaries.  On
the first leg,
$\|g_{2j+1}-g_{2j}\|=\|y_{2j}\|=O(\epsilon_j^2)$, so
\[
\|g_k\|\longrightarrow G_\infty>0.
\]
All $H_k$ remain positive definite because $s_k^Ty_k=s_k^TA_ks_k>0$ and
DFP preserves positive definiteness.  This establishes the construction in
dimension two.

For $n>2$, use
\[
\widehat f(z,w)=f(z)+\frac12\|w\|^2,
\qquad
\widehat H_0=H_0\oplus I,
\qquad w_0=0.
\]
The entire iteration remains in the first two coordinates, so all identities
and inequalities are unchanged.  This completes the proof of
Theorem~\ref{thm:main}.

\subsection{Relation to known sufficient conditions}
\label{sec:known-results}

The orbit lies outside the additional hypotheses in the known conditional
convergence theorems.  The exact expansions in
Appendix~\ref{app:algebra}, normalized by the cycle amplitude $G_j$, give
\begin{align*}
\norm{g_{2j}}/G_j
 &=1+2\epsilon_j^4+O(\epsilon_j^6),\\
\norm{g_{2j+1}}/G_j
 &=1-2\epsilon_j^3-2\epsilon_j^4+O(\epsilon_j^6),\\
\norm{g_{2j+2}}/G_j
 &=1-\frac92\epsilon_j^4+O(\epsilon_j^6).
\end{align*}
Thus the first leg decreases the gradient norm by order $\epsilon_j^3$ and
the second increases it by order $\epsilon_j^3$; eventual monotonicity fails.
Moreover,
\[
\norm{s_{2j}}/G_j=2\epsilon_j^2+o(\epsilon_j^2),
\qquad
\norm{s_{2j+1}}/G_j=\epsilon_j^2+o(\epsilon_j^2).
\]
Since $G_j\to G_\infty>0$ and
$\sum_j\epsilon_j^2=\infty$, the total step length
$\sum_k\norm{s_k}$ diverges.  Hence neither the eventual gradient-norm
condition nor the summable-step condition in Yuan \cite{Yuan1995} applies.

For the final function, $m=1/2$ and $M=3/2$.  Xu's theorem
\cite{Xu1997} assumes, in addition to an eventual monotonicity condition on a
trajectory-dependent scalar, that $c_2<(m/M)^3$.  Here
\[
c_2=\frac34>\left(\frac{1/2}{3/2}\right)^3=\frac1{27},
\]
so its parameter hypothesis is violated.  Powell's theorem
\cite{Powell1971} uses exact line search, whereas our two phases satisfy
$g_{k+1}^Ts_k=-(1-\tau_k)q_k\ne0$.  Finally, the global result of Byrd,
Nocedal, and Yuan \cite{ByrdNocedalYuan1987} excludes the DFP endpoint.
Consequently the counterexample does not contradict any of these positive
results; it identifies why their extra hypotheses matter.

\subsection{Verification of the assumptions}

The construction uses only the following fixed data:
\begin{itemize}
\item the unmodified inverse DFP formula;
\item one fixed symmetric positive definite $H_0$ obtained from
  \eqref{eq:canonical-state} on the invariant curve;
\item the standard weak Wolfe constants $c_1=1/4$ and $c_2=3/4$;
\item a globally $C^2$ objective with the fixed bounds
  $\frac12I\preceq\nabla^2f\preceq\frac32I$.
\end{itemize}
No exact line search, strong Wolfe condition, damping, restart, projection,
or bounded-conditioning assumption is used.  The matrices in
\eqref{eq:two-controls} are auxiliary secant matrices; the
actual function is \eqref{eq:objective}.  Search segments may cross bump
supports, which is immaterial because weak Wolfe uses endpoint values and
endpoint directional derivatives, while the fundamental theorem of calculus
recovers the exact secant from the final $C^2$ function.

\section{Conclusion}

The construction gives a globally uniformly convex $C^2$ objective for which
the unmodified classical DFP method satisfies the standard weak Wolfe
conditions at every step while its gradient norms converge to a positive
constant.  Thus weak Wolfe line search and uniform convexity alone do not
guarantee global convergence of DFP.  The counterexample isolates a singular
two-phase metric drift that is excluded by exact-search orthogonality and by
the additional trajectory assumptions in earlier positive results; sharper
Wolfe-parameter ranges and smoother realizations remain separate questions.

\section*{Statements and Declarations}

\textbf{Funding.}
The authors will insert the complete funding statement, or state that no
funding was received, after confirming the submission metadata.

\textbf{Competing interests.}
The authors will insert their confirmed financial and non-financial competing
interests statement before submission.

\textbf{Availability of data and materials.}
No datasets were generated or analyzed.  A deterministic symbolic
verification program reproducing the coefficients in
Appendix~\ref{app:algebra} is available from the corresponding author and is
included as Additional Material for referees.  The proof in the article does
not depend on executing that program.

\textbf{Use of large-language-model tools.}
Large-language-model tools were used to assist with symbolic-computation
code, adversarial proof checking, literature organization, and language
editing.  The human authors checked and approved every formula, citation,
proof, and conclusion in the submitted version and take full responsibility
for the content.

\appendix

\section{Algebraic verification of the two-leg map}
\label{app:algebra}

This appendix makes the calculation behind Lemmas~\ref{lem:slow-curve} and
\ref{lem:physical-drift} directly auditable from the printed DFP formula.
The symbolic verification program supplied to referees evaluates the same
rational identities; it is a reproducibility aid, not a premise of the proof.

\subsection{One step in entries and spectral coordinates}

Work in an eigenframe and write
\[
H=\begin{pmatrix}\ell&0\\0&\eta\end{pmatrix},\qquad
g=\binom{g_1}{g_2},\qquad
A=\begin{pmatrix}a&c\\c&d\end{pmatrix}.
\]
Set
\[
v_1=\ell g_1,\quad v_2=\eta g_2,\quad
w_1=av_1+cv_2,\quad w_2=cv_1+dv_2,
\]
and
\[
\delta=g_1v_1+g_2v_2,\qquad
\beta=v_1w_1+v_2w_2,
\qquad \gamma=\ell w_1^2+\eta w_2^2.
\]
Proposition~\ref{prop:one-step} gives, entry by entry,
\begin{align}
(H_+)_{11}&=\ell-\frac{\ell^2w_1^2}{\gamma}
                    +\frac{v_1^2}{\beta},
&
(H_+)_{12}&=-\frac{\ell\eta w_1w_2}{\gamma}
                    +\frac{v_1v_2}{\beta},
\label{eq:appendix-H-entries}\\
(H_+)_{22}&=\eta-\frac{\eta^2w_2^2}{\gamma}
                    +\frac{v_2^2}{\beta},
&
(g_+)_i&=g_i-\tau\frac{\delta}{\beta}w_i.
\nonumber
\end{align}
These identities also give
\begin{equation}\label{eq:dfp-determinant}
\det H_+=\det H\,\frac{s^Ty}{y^THy},
\end{equation}
either by a two-by-two determinant expansion or by the matrix determinant
lemma.

For a symmetric matrix
$K=\left(\begin{smallmatrix}a_0&b_0\\b_0&d_0\end{smallmatrix}\right)$,
put
\[
\Delta=\sqrt{(d_0-a_0)^2+4b_0^2},\qquad
\lambda_\pm=\frac{a_0+d_0\pm\Delta}{2}.
\]
Near the limiting matrix $\operatorname{diag}(0,1)$, an analytic low
eigenvector is
\[
u_-=
\frac{(d_0-\lambda_-,-b_0)^T}
{\sqrt{(d_0-\lambda_-)^2+b_0^2}},
\qquad u_+=u_-^\perp.
\]
Reverse both signs if needed so that $\gamma_-=u_-^Tg>0$, and write
$\gamma_+=u_+^Tg$.  Comparing with \eqref{eq:canonical-state} recovers the
canonical variables by the exact formulas
\begin{equation}\label{eq:canonical-recovery}
G=\gamma_-,\qquad h=\lambda_+,qquad
r=\frac{\lambda_-\gamma_-}{\lambda_+\gamma_+},
\qquad
p=\frac{\lambda_+\gamma_+^2}{\lambda_-\gamma_-^2}.
\end{equation}
Equations \eqref{eq:appendix-H-entries} and
\eqref{eq:canonical-recovery} are the complete one-leg calculation used on
both phases.

\subsection{Two-leg expansion and the invariant graph}

To expose the mixed orders, replace $\epsilon$ in the two controls by an
independent variable $b$, leave the canonical radius $r$ independent, and set
\[
p=2+Pbr,\qquad h=1+Jbr.
\]
Applying \eqref{eq:appendix-H-entries}, diagonalizing with the preceding
quadratic formula, and recovering the canonical variables with
\eqref{eq:canonical-recovery} gives
\begin{align}
r_+-r
 &=\frac{b(6J+5P-300)}{18}r^2+O(r^3),
\label{eq:appendix-r}\\
p_+-2
 &=\frac{b(6J-P+348)}9r+O(r^2),
\label{eq:appendix-p}\\
h_+-1&=8br+O(r^2),
\label{eq:appendix-h}\\
\frac{G_+}{G}-1
 &=\frac{b^2(24J-4P+384)-117}{18}r^2+O(r^3),
\label{eq:appendix-G}\\
\phi_+-\phi&=-3r+O(r^2),
\label{eq:appendix-phi}\\
e^T(C_+-C)
 &=-\frac{2b^2(6J-P+96)}9Gr^2+O(G|b|r^3).
\label{eq:appendix-C}
\end{align}
The remainders are uniform for bounded $(b,P,J)$ in the fixed Taylor
neighborhood.  In particular, \eqref{eq:appendix-C} is a joint estimate; only
after setting $b=\epsilon$ and $r=\epsilon^2$ is its remainder
$o(G\epsilon^6)$.

For the graph ansatz
\[
p=2+P_3\epsilon^3+P_4\epsilon^4+O(\epsilon^5),\qquad
h=1+H_3\epsilon^3+H_4\epsilon^4+O(\epsilon^5),
\]
the coefficients of $p_+-p(\epsilon_+)$ and
$h_+-h(\epsilon_+)$ are as follows:
\begin{center}
\begin{tabular}{c@{\qquad}c@{\qquad}c}
\toprule
order & $p$-component & $h$-component\\
\midrule
$\epsilon^3$ & $\frac23H_3-\frac{10}{9}P_3+\frac{116}{3}$
              & $8-H_3$\\[2pt]
$\epsilon^4$ & $\frac23H_4-\frac{10}{9}P_4-2$
              & $-H_4$\\
\bottomrule
\end{tabular}
\end{center}
Setting these four entries to zero yields
$(P_3,H_3,P_4,H_4)=(198/5,8,-9/5,0)$, exactly as in
\eqref{eq:graph-coefficient-system}.

\subsection{Specialized coefficients on the slow curve}

Substitution of those four graph coefficients into the exact two-leg formulas
gives the following expansions, all normalized by the cycle-boundary
amplitude $G_j$ where appropriate:
\begin{align*}
\epsilon_+
 &=\epsilon-\frac32\epsilon^4+\frac54\epsilon^5+O(\epsilon^6),\\
\frac{G_+}{G}
 &=1-\frac{13}{2}\epsilon^4+\frac{116}{5}\epsilon^6
   -\frac{976}{5}\epsilon^7+O(\epsilon^8),\\
\phi_+-\phi
 &=-3\epsilon^2-\frac{196}{5}\epsilon^5
   +\frac{28}{5}\epsilon^6+O(\epsilon^7),\\
R_j^T(C_{2j+2}-C_{2j})/G_j
 &=\binom{-\frac{116}{5}\epsilon^6+\frac{38}{5}\epsilon^7
           +O(\epsilon^8)}
          {-\frac{508}{5}\epsilon^8+O(\epsilon^9)}.
\end{align*}
The two endpoint angle increments are
\begin{align*}
-2\epsilon^2-\frac{122}{5}\epsilon^5
 +\frac{88}{15}\epsilon^6+O(\epsilon^7),\qquad
-\epsilon^2-\frac{104}{5}\epsilon^5
 +\frac{71}{15}\epsilon^6+O(\epsilon^7).
\end{align*}
The corresponding step and gradient norms are
\begin{align*}
\frac{\norm{s_{2j}}}{G_j}
 &=2\epsilon^2+\frac{112}{5}\epsilon^5
   -\frac{11}{5}\epsilon^6+O(\epsilon^7),\\
\frac{\norm{s_{2j+1}}}{G_j}
 &=\epsilon^2+\frac{114}{5}\epsilon^5
   -\frac{49}{10}\epsilon^6+O(\epsilon^7),\\
\frac{\norm{g_{2j}}}{G_j}
 &=1+2\epsilon^4+O(\epsilon^6),\\
\frac{\norm{g_{2j+1}}}{G_j}
 &=1-2\epsilon^3-2\epsilon^4-\frac{112}{5}\epsilon^6
   +O(\epsilon^7),\\
\frac{\norm{g_{2j+2}}}{G_j}
 &=1-\frac92\epsilon^4+\frac{116}{5}\epsilon^6
   +O(\epsilon^7).
\end{align*}
Finally, direct substitution in \eqref{eq:abstract-step} gives the exact
ratios $s^Ty/(-g^Ts)=2/3$ and $1/3$.  Thus every coefficient used in the
slow dynamics, separation argument, known-results comparison, and Wolfe
verification follows from the entry formulas printed in this appendix.

\begin{thebibliography}{11}

\bibitem{ByrdNocedalYuan1987}
Byrd, R.H., Nocedal, J., Yuan, Y.-X.:
Global convergence of a class of quasi-Newton methods on convex problems.
SIAM J. Numer. Anal. 24(5), 1171--1190 (1987).
\url{https://doi.org/10.1137/0724077}

\bibitem{Davidon1991}
Davidon, W.C.:
Variable metric method for minimization.
SIAM J. Optim. 1(1), 1--17 (1991).
\url{https://doi.org/10.1137/0801001}

\bibitem{FletcherPowell1963}
Fletcher, R., Powell, M.J.D.:
A rapidly convergent descent method for minimization.
Comput. J. 6(2), 163--168 (1963).
\url{https://doi.org/10.1093/comjnl/6.2.163}

\bibitem{Fletcher1994}
Fletcher, R.:
An overview of unconstrained optimization.
In: Spedicato, E. (ed.) Algorithms for Continuous Optimization,
pp. 109--143.
Springer, Dordrecht (1994).
\url{https://doi.org/10.1007/978-94-009-0369-2_5}

\bibitem{KarydasSchinas1992}
Karydas, N., Schinas, J.:
The center manifold theorem for a discrete system.
Appl. Anal. 44(3--4), 267--284 (1992).
\url{https://doi.org/10.1080/00036819208840083}

\bibitem{Nocedal1992}
Nocedal, J.:
Theory of algorithms for unconstrained optimization.
Acta Numer. 1, 199--242 (1992).
\url{https://doi.org/10.1017/S0962492900002270}

\bibitem{Powell1971}
Powell, M.J.D.:
On the convergence of the variable metric algorithm.
J. Inst. Math. Appl. 7(1), 21--36 (1971).
\url{https://doi.org/10.1093/imamat/7.1.21}

\bibitem{Wolfe1969}
Wolfe, P.:
Convergence conditions for ascent methods.
SIAM Rev. 11(2), 226--235 (1969).
\url{https://doi.org/10.1137/1011036}

\bibitem{Xu1997}
Xu, D.-C.:
Global convergence analysis of DFP method.
Math. Numer. Sin. 19(3), 287--292 (1997).
\url{https://doi.org/10.12286/jssx.1997.3.287}

\bibitem{Yuan1995}
Yuan, Y.-X.:
Convergence of DFP algorithm.
Sci. China Ser. A 38(11), 1281--1294 (1995).

\bibitem{YuanZhouPham2026}
Yuan, G., Zhou, P., Pham, H.:
A projection DFP quasi-Newton algorithm and its applications in Muskingum
model and machine learning.
Numer. Algorithms (2026).
\url{https://doi.org/10.1007/s11075-025-02284-6}

\end{thebibliography}

\end{document}
